% Generated draft: requires mathematical and source review.
\documentclass[11pt]{article}
\usepackage[margin=25mm]{geometry}
\usepackage{fontspec}
\setmainfont{Latin Modern Roman}
\usepackage{amsmath,amssymb,mathtools,unicode-math}
\setmathfont{Latin Modern Math}
\usepackage{enumitem,needspace,xurl,hyperref}
\setlength{\parindent}{0pt}
\setlength{\parskip}{5pt}
\setlength{\emergencystretch}{3em}
\newcommand{\N}{\mathbb{N}}
\newcommand{\Z}{\mathbb{Z}}
\newcommand{\Q}{\mathbb{Q}}
\newcommand{\R}{\mathbb{R}}
\newcommand{\C}{\mathbb{C}}
\newcommand{\F}{\mathbb{F}}
\newcommand{\A}{\mathbb{A}}
\newcommand{\PP}{\mathbb{P}}
\newcommand{\E}{\mathbb{E}}
\newcommand{\eps}{\varepsilon}
\newcommand{\dd}{\,\mathrm{d}}
\begin{document}
\section*{UnsolvedMath ID: 30002322}
\textbf{Source problem code:} catalogue entry
\subsection{Definitions and mathematical statement}
Let \(G=(V,E)\) be a finite simple undirected graph, and write \(n=|V|\). The graph is triangle-free if no three distinct vertices are pairwise adjacent. Disjoint subsets \(A,B\subseteq V\) are anticomplete if \(\{a,b\}\notin E\) for every \(a\in A\) and \(b\in B\); no restriction is imposed on edges within either subset.

Under the lower-bound interpretation of the supplied phrase ``of orders \(cn^c\) and \(cn\)'', the question is whether there exists a real constant \(c>0\), independent of \(n\) and \(G\), such that every triangle-free graph \(G\) with \(n>2\) admits disjoint anticomplete subsets \(A,B\subseteq V\) satisfying
\[
|A|\geq cn^c\qquad\text{and}\qquad |B|\geq cn.
\]
Because cardinalities are integers, these inequalities are equivalent to \(|A|\geq\lceil cn^c\rceil\) and \(|B|\geq\lceil cn\rceil\). The same constant \(c\) occurs in both coefficients and in the exponent. The source does not explicitly specify a rounding convention or say ``at least''; thus this lower-bound interpretation requires confirmation. The finite simple undirected graph convention is also made explicit here rather than stated in the supplied wording.
\subsection{Short English statement}
Does one universal constant \(c>0\) guarantee that every triangle-free graph on \(n>2\) vertices contains two disjoint vertex sets, one of size at least \(cn^c\) and the other of size at least \(cn\), with no edges between them? This uses an explicitly qualified lower-bound interpretation of the source's size wording.
\subsection{Sources}
\begin{itemize}
\item Local UnsolvedMath registry, permanent ID \(30002322\), title ``Large Anticomplete Sets in Triangle-Free Graphs'', source problem code \texttt{catalogue entry}, in the Oberwolfach Reports Open Problems collection. The record supplies proposed year 2013 and no individual proposer.
\item Supplied report citation: \emph{Combinatorics and Probability (2014)}, Oberwolfach Reports, DOI \texttt{10.4171/owr/2013/18}, pp.\ 1087--1152. \url{https://doi.org/10.4171/owr/2013/18}. The supplied background describes these reports as documenting mathematical workshops and their open-problem discussions. The citation's 2014 label and its 2013 DOI component are retained as supplied.
\item Supplied literature reference: Chudnovsky--Fox--Scott--Seymour--Spirkl, ``anticomplete sets''. \url{https://arxiv.org/abs/1810.00058}.
\item Supplied literature reference: ``Oberwolfach combinatorics report''. \url{https://doi.org/10.4171/owr/2013/18}.
\item Literature assessment dated exactly \texttt{2026-08-21}: the supplied curation labels the problem open and describes it as the triangle case of the Conlon--Fox--Sudakov anticomplete-sets conjecture. This is supplied dated curation derived from the OWR material, not an independently verified literature assessment; the background explicitly requests independent verification of time-sensitive status and cited claims.
\item The registry labels the statement \texttt{corrected\_verified} and reports that its mathematical content was checked against the official Oberwolfach report, with notation, context, and question form normalized without changing the source-level claim. This verification description is supplied metadata, not verification performed for this draft.
\end{itemize}
\end{document}
